\chapter{Introduction}
\label{ch:intro}

\section{Why procurement disputes are worth modelling}
\label{sec:context}

Public bodies in the United Kingdom buy goods, works and services worth roughly
\pounds300 billion a year \cite{cabinetoffice2020}, and the rules governing how
they do it were rewritten in 2023. The Procurement Act 2023 \cite{pa2023} replaced a patchwork of
EU-derived regulations with a single regime built around four statutory
objectives in section 12: value for money, maximising public benefit, sharing
information for the purpose of transparency, and acting with integrity. It also
tightened the machinery around losing bidders. Section 50 requires a contracting
authority to give each unsuccessful supplier an assessment summary explaining how
its tender was evaluated; section 51 opens a standstill period of eight working
days beginning with the day the contract award notice is published, during which
the contract cannot be signed. A supplier who believes the evaluation went wrong
has that window to raise a concern, and thirty days from the date it knew or
ought to have known of the breach to start proceedings in the Technology and
Construction Court.

Those deadlines are the reason this problem is interesting rather than merely
legal. Figure~\ref{fig:lifecycle} shows the shape of the timeline. Almost all of
the leverage sits in the first fortnight, before anybody has instructed a
solicitor, and almost all of the information asymmetry sits there too: the
authority knows what its evaluation panel actually did, and the losing bidder
knows only what the assessment summary chose to tell it. Once the automatic
suspension bites and the case reaches the TCC, both parties are committed to a
process that in this project's own case corpus took between four months and two
years to conclude.

\begin{figure}[htbp]
\centering
\begin{tikzpicture}[font=\footnotesize,>=Latex]
  % timeline spine
  \draw[line width=1.1pt,f21ink] (0,0) -- (13.2,0);
  \foreach \x in {0,2.6,5.2,8.4,11.6} \draw[f21ink,line width=1.1pt] (\x,-0.13) -- (\x,0.13);

  % labels below
  \node[anchor=north,text width=2.3cm,align=center] at (0,-0.25)
    {Evaluation\\ complete};
  \node[anchor=north,text width=2.4cm,align=center] at (2.6,-0.25)
    {Assessment\\ summary sent\\ \textit{(s.50)}};
  \node[anchor=north,text width=2.4cm,align=center] at (5.2,-0.25)
    {Standstill ends\\ 8 working days\\ \textit{(s.51)}};
  \node[anchor=north,text width=2.5cm,align=center] at (8.4,-0.25)
    {Claim issued\\ 30-day limit\\ suspension bites};
  \node[anchor=north,text width=2.5cm,align=center] at (11.6,-0.25)
    {TCC judgment\\ 4 months--2 years};

  % artefact bands above
  \begin{scope}[on background layer]
    \fill[f21green!12,rounded corners=3pt] (-0.5,0.42) rectangle (3.9,1.5);
    \fill[f21amber!12,rounded corners=3pt] (4.3,0.42) rectangle (12.6,1.5);
  \end{scope}
  \node[anchor=west,text width=4.2cm,align=left,f21green] at (-0.45,0.96)
    {\textbf{Risk screen} (\S\ref{sec:risk-screen})\\ acts before a dispute exists};
  \node[anchor=west,text width=8.0cm,align=left,f21amber] at (4.4,0.96)
    {\textbf{Negotiation simulator} (\S\ref{sec:orchestration}) acts on a dispute\\ that has already been raised};
\end{tikzpicture}
\caption[The procurement dispute lifecycle]{The procurement dispute lifecycle
under the Procurement Act 2023, and where this project's two artefacts sit on it.
The statutory windows are short; the litigation that follows is not. The horizontal
axis is ordinal, not to scale.}
\label{fig:lifecycle}
\end{figure}

This project was carried out in partnership with Fusion21, a social enterprise
that runs procurement frameworks for housing providers, local authorities and
NHS bodies, and which sees the same disputes recur across its member base.
Fusion21 supplied four internal documents, referred to throughout as Doc 1--4,
covering the interests and drivers of each party to a dispute, the escalation
ladder from informal engagement through to litigation, the negotiation theory
they apply, and three worked examples. Doc 4 makes an observation that quietly
constrains the whole project: litigated cases are not typical, because the
overwhelming majority of procurement disputes settle before proceedings or
before trial. Every judgment used in this evaluation is therefore a survivor of a
filtering process, and any claim built on that corpus inherits the bias.

\section{Subject area and related work}
\label{sec:related}

The technical subject here is the use of several large language model agents,
each given a role, a set of interests and a private fallback position, to
simulate a structured multi-party interaction. Three strands of prior work bear
on it directly.

\subsection*{Multi-agent LLM frameworks}

Since instruction-following models became capable of few-shot task specification
\cite{brown2020language} and of exposing intermediate reasoning on request
\cite{wei2022chain}, the prompt has functioned as a programming surface. The
observation that a single model can then be instantiated several times
with different role prompts, and that the resulting agents will produce something
resembling a division of labour, is the basis of CAMEL \cite{li2023camel}, which
paired an assistant and a user agent under an inception-prompting scheme, and of
Generative Agents \cite{park2023generative}, which added memory, reflection and
planning to twenty-five simulated townspeople. AutoGen \cite{wu2023autogen} and
MetaGPT \cite{hong2024metagpt} pushed the same idea toward engineering
workflows, MetaGPT in particular arguing that encoding standard operating
procedures into the prompts is what stops multi-agent conversation from
degenerating. That argument turns out to matter here: the Court agent's
verification protocol in Section~\ref{sec:court-design} is precisely such an
encoded procedure --- closer to ReAct's interleaving of reasoning with a defined
action space \cite{yao2023react} than to open-ended chat --- and
Chapter~\ref{ch:evaluation} shows what happens on the case
shapes it was not written for. Orchestration in this project uses LangGraph
\cite{langgraph}, which models the interaction as an explicit state graph rather
than as free-form conversation, a choice justified in Section~\ref{sec:orchestration}.

\subsection*{LLMs as negotiators}

Negotiation is a natural testbed because it has structure, an outcome, and a
literature. Fu et al.\ \cite{fu2023negotiation} had two models bargain over a
balloon and improve through self-play with AI feedback. NegotiationArena
\cite{bianchi2024negotiation} built a benchmark of resource-allocation,
buy--sell and ultimatum games and found LLM negotiators susceptible to
manipulation tactics such as feigned desperation. Abdelnabi et al.\
\cite{abdelnabi2024negotiation} moved closer to the setting here, constructing
multi-party negotiations with conflicting private goals and scoring agents on
whether they reached a deal that satisfied their own constraints. The common
thread is that these environments are self-contained: the payoff is defined by
the game, so correctness is decidable. Procurement disputes are not like that.
The governing framework is external, contested, and expressed in statute and case
law, and the correct outcome is whatever a court would have held, which is
usually unknown.

The normative framework the agents reason with comes from Fisher and Ury's
\emph{Getting to Yes} \cite{fisher2011getting} --- separate interests from
positions, know your BATNA --- and from Raiffa's treatment of asymmetric
negotiation \cite{raiffa1982art}. Doc 3 imposes a hard constraint from that
literature which shapes the system's outcome vocabulary: in award-stage
procurement disputes you cannot split the difference, because the remedies
available are fixed by statute and re-running an award is not a bargaining chip.

\subsection*{Legal reasoning and the fabrication problem}

Language models perform respectably on legal benchmarks. GPT-4 passed the Uniform
Bar Examination \cite{katz2024gpt4}, and LegalBench \cite{guha2023legalbench}
assembled 162 tasks showing real but uneven legal reasoning ability. The
countervailing evidence is more relevant to this project. Dahl et al.\
\cite{dahl2024legal} profiled legal hallucination across models and found error
rates between 58\% and 82\% on verifiable questions about federal court cases,
with models systematically failing to recognise when they did not know. Magesh
et al.\ \cite{magesh2024hallucination} tested commercial retrieval-augmented
legal research tools and still found hallucination rates above 17\%. The general
survey position \cite{ji2023survey} is that fluent, confident, wrong output is
the default failure mode, not an edge case.

That is the whole difficulty. A system that produces a plausible compliance
finding on a procurement dispute is useless --- worse than useless, because it is
persuasive --- unless something guarantees that the numbers and provisions it cites
actually came from the case. Retrieval-augmented generation \cite{lewis2020rag}
is the standard mitigation and was considered here; Section~\ref{sec:alternatives}
explains why it was scoped out and what was done instead.

\section{The gap this project addresses}
\label{sec:gap}

Existing multi-agent negotiation work evaluates against a payoff function the
environment defines. Existing legal-LLM work evaluates against benchmark answers a
human labelled. Neither situation holds for a simulated procurement dispute, where
the artefact produces a legal judgment on facts that a real court either decided
differently, decided on a different question, or never decided at all. The
practical question this project ended up organising itself around is therefore
not \emph{can agents negotiate a procurement dispute}, which turns out to be
comparatively easy, but \emph{how do you evaluate what they produce, and how do
you stop them inventing the evidence they reason from}.

\section{Aim, research questions and objectives}
\label{sec:objectives}

\textbf{Aim.} To build and evaluate a multi-agent LLM system that simulates
procurement dispute resolution under the Procurement Act 2023, modelling the
interests, constraints and fallback positions of the contracting authority, the
aggrieved bidder and the court, and to establish an evaluation methodology
appropriate to a domain where ground truth is scarce and fabrication is the
principal risk.

The research questions carried forward from the project plan, lightly sharpened
by what the work uncovered, are:

\begin{description}
  \item[RQ1] Can prompt engineering and structured context injection, without
    fine-tuning, produce agents that reliably represent the distinct interests,
    legal constraints and BATNA of each party?
  \item[RQ2] What negotiation protocol allows autonomous agents to reach
    legally grounded resolutions, and what component of that protocol is actually
    load-bearing?
  \item[RQ3] Can an LLM court agent be constrained to verify arithmetic when the
    case supplies one and refuse to invent one when it does not, and does that
    constraint hold across dispute shapes it was not designed against?
  \item[RQ4] How should outcomes be evaluated for reliability, fabrication and
    alignment with real dispositions, given that most real disputes settle
    confidentially and never produce a ground-truth label?
\end{description}

These decompose into six objectives, which double as the structure of
Chapter~\ref{ch:methodology}: (O1) specify three negotiating agents plus a
non-negotiating summariser against Fusion21's own interest taxonomies; (O2)
implement a deterministic, inspectable orchestration protocol with structured,
schema-validated output; (O3) design a court agent that verifies conditionally
rather than unconditionally; (O4) build an ingestion path from real documents to
structured scenarios; (O5) assemble a verified corpus of real UK judgments and
evaluate against it with explicit baselines and ablations; and (O6) deliver
something Fusion21 can actually use, which turned out to mean a pre-award risk
screen rather than the simulator.

\section{Contributions}
\label{sec:contributions}

\begin{enumerate}
  \item \textbf{A conditional verification protocol for an LLM adjudicator}
    (Section~\ref{sec:court-design}) that gates exact arithmetic on the presence
    of both a formula and every input it requires, and forbids invented values
    even when hedged as illustrative. Four prompt revisions are reported, each
    with the measured failure mode that motivated the next.
  \item \textbf{An evaluation methodology for a scarce-ground-truth domain}
    (Chapter~\ref{ch:evaluation}), combining agreement with real dispositions,
    three independent fabrication screens, structural-compliance instrumentation,
    a controlled prompt ablation with significance testing, and three baselines.
  \item \textbf{A negative result that constrains the architecture claim}
    (Section~\ref{sec:baselines}): a naive zero-shot prompt matches the full
    pipeline's directional accuracy, while removing the court agent deadlocks
    every run. The architecture earns itself on resolution mechanics and process,
    not on accuracy.
  \item \textbf{A taxonomy of three distinct failure modes} --- numeric
    fabrication, narrative-premise fabrication, and legal-judgment divergence ---
    each evidenced on named cases (Section~\ref{sec:fabrication}), together with
    the finding that hardening the adjudicator against the first does nothing
    about the second.
  \item \textbf{A verified corpus of 23 real UK procurement judgments} rendered
    as structured scenarios, with dispositions checked individually, released with
    the project.
  \item \textbf{A pre-award challenge risk screen} (Section~\ref{sec:risk-screen})
    traced rule by rule to Fusion21's own practical-considerations tables, and the
    only component of this system with an observable ground truth in principle.
\end{enumerate}

\section{Report structure}

Chapter~\ref{ch:methodology} describes the system and justifies the design
choices against the alternatives that were considered and rejected.
Chapter~\ref{ch:evaluation} presents the evaluation, including results that do
not favour the artefact. Chapter~\ref{ch:conclusion} returns to the objectives,
reflects on the process, and sets out what should be built next.
